\subsection{Metodología Methadis}
Es una metodología para diseñar un sistema hipermedia adaptativo para el aprendizaje basado en estilos de aprendizaje y estilos cognitivos.

El diseño se divide en tres fases \cite{Ferraro2006}:

\textbf{Fase 1: Definición de la unidad de aprendizaje}
\begin{itemize}
    \item Definición del objetivo de aprendizaje: Se establece lo que se quiere que el estudiante aprenda al final de la unidad y se define: El dominio de aprendizaje (se determina si el aprendizaje corresponde al dominio cognitivo, afectivo o motor), la taxonomía del tipo de aprendizaje y tipo de aprendizaje específico.

    \item Definir el Modelo del Dominio: a partir del objetivo de aprendizaje, se definen los contenidos y su estructura para el modelo del dominio. Se debe hacer una lista de conceptos o actividades, dividir cada uno de estos conceptos o actividades de aprendizaje en nodos o páginas, estas se descomponen en actividades específicas o conceptos más pequeños y se establecen las relaciones entre ellas y, por último, se representa en forma esquematizada.

\end{itemize} \ \\
\textbf{Fase 2: Selección de las teorías de estilos de aprendizaje y estilos cognitivos}
\begin{itemize}
    \item Selección de las teorías de estilos de aprendizaje: Elección de las teorías de estilos de aprendizaje y estilos cognitivos que serán utilizadas para clasificar a los estudiantes en diferentes categorías de usuarios. Esta elección se debe hacer mediante los siguientes criterios: 
\begin{itemize}
    \item ¿Es apropiado para el contexto de aprendizaje?
    \item Justificación teórica y empírica ¿qué tipo de personas o muestras se les puede aplicar el instrumento de medición?
    \item Relación de cada estilo con las correspondientes estrategias instruccionales.
    \item Disponibilidad de instrumentos de medición.
    \item Cantidad de ítems del instrumento de medición.
    \item Costo del instrumento de medición. 
\end{itemize}
\item Estos criterios se aplican para selección a las teorías de estilos de aprendizaje, los cuales se adaptará el sistema.
\end{itemize} \ \\
\textbf{Fase 3: Adaptación para el aprendizaje}
\begin{itemize}
\item Definir el modelo del usuario:  Al seleccionar las teorías de estilos de aprendizaje y de estilos cognitivos, se debe aplicar los instrumentos de medición de las categorías de estos para así determinar el estilo de aprendizaje y cognitivo del estudiante.

\item Definir el Modelo de adaptación: Se seleccionan las estrategias de enseñanza a cada estilo de aprendizaje y cada estilo cognitivo según sus preferencias. Cabe destacar que la presentación de contenidos se ajusta a cada uno de los estilos de aprendizaje y las opciones de navegación se ajustas a cada una de las categorías de estilos cognitivos.
\end{itemize}